\section{总结与延伸}

本讲从预训练 scaling laws 出发，解释了大模型能力为何曾长期依赖参数、数据和训练计算的共同扩展；随后通过 zero/few-shot 与 chain-of-thought 展示上下文如何激活模型内部潜在计算；通过 instruction tuning 与 RLHF 说明后训练如何塑造行为；再用 Large Language Monkeys 和 reasoning model 引出推理时计算；最后把这些能力放入带工具和反馈的 Agent 系统。

最值得保留的三个判断是：

\begin{enumerate}
    \item \textbf{能力不只存在于权重，也存在于可调用的计算过程。} 同一模型在不同采样、搜索和推理预算下可呈现显著不同的成功率。
    \item \textbf{生成能力和验证能力必须分开测量。} 候选中“有正确答案”不等于系统“能交付正确答案”。
    \item \textbf{Agent 的价值来自闭环。} 只有真实执行、反馈消费和状态更新，才能把语言模型升级为可持续完成任务的系统。
\end{enumerate}

下一讲将专门讨论 test-time compute scaling：重复采样为何呈幂律收益、并行采样与串行修订如何比较、reward model 如何指导搜索，以及何时额外预训练仍优于额外推理计算。

\end{document}
